% Invented example for the complete-paper skill (parts/supplement.md).
% The physics, the numbers and the references are made up; only the use of the supplement is the point.
\documentclass[twocolumn,10pt]{article}
\usepackage[margin=18mm]{geometry}
% One preamble for main.tex and supplement.tex: the same packages, the same notation, the same macros.
% Invented example for the complete-paper skill.
\usepackage{amsmath}
\usepackage{graphicx}
\usepackage{xr}
\newcommand{\jc}{j_{\mathrm{c}}}        % threshold current density
\newcommand{\jzero}{j_{0}}              % threshold at 90 degrees of tilt
\newcommand{\mz}{m_z}                   % out-of-plane component of the octupole moment
\newcommand{\tsh}{\theta_{\mathrm{SH}}} % spin Hall angle
\newcommand{\Ku}{K}                     % anisotropy energy density
\newcommand{\un}[1]{\,\mathrm{#1}}      % units: upright, thin space

\externaldocument[S-]{supplement}

\title{Field-free switching of a tilted kagome antiferromagnet}
\author{A. Author \and B. Author}
\date{}

\begin{document}
\maketitle

\begin{abstract}
Electrical switching of noncollinear antiferromagnets has so far required an in-plane magnetic field, which
blocks their use in dense memory. Here we show by micromagnetic simulation that tilting the kagome plane of
Mn$_3$Sn away from the film normal removes this requirement. A current alone then reverses the octupole
moment, above a threshold that falls as the inverse sine of the tilt angle. The threshold agrees with an
analytical balance of spin-orbit torque and anisotropy. Tilt is therefore a design parameter for field-free
antiferromagnetic memory.
\end{abstract}

\section{Introduction}

Noncollinear antiferromagnets combine the fast dynamics of antiferromagnets with a readout signal as large
as that of ferromagnets. In Mn$_3$Sn the signal comes from a cluster octupole moment, which a spin-orbit
torque can rotate. Every switching experiment reported so far applies a magnetic field along the current to
make the reversal deterministic~\cite{A,B}. A field is acceptable in a test structure and unacceptable in a
memory array, where neighbouring cells would share it.

One way to break the symmetry without a field is to break it in the crystal. Films of Mn$_3$Sn grow with
the kagome plane tilted from the film normal when the substrate is miscut~\cite{C}. Building memory from such films
needs two answers, whether the tilt makes the current-induced reversal deterministic and at what cost in
current. We answer both questions with simulations of a Mn$_3$Sn/W bilayer and with an analytical model.

\section{Model}

Figure~\ref{fig:device} shows the geometry. A current density $j$ flows along $x$ in the W layer and injects
a spin current polarized along $y$ into the Mn$_3$Sn film of thickness $d$. The kagome plane is tilted by
the angle $\theta$ about the $y$ axis. We integrate the coupled equations of motion of the three sublattice
moments at zero temperature unless stated otherwise.

\begin{figure}
  \centering
  \includegraphics[width=0.9\columnwidth]{example-image-a}
  \caption{Geometry of the bilayer. The kagome plane (shaded) is tilted by the angle $\theta$ about the $y$
  axis, and the current density $j$ flows along $x$. Not to scale.}
  \label{fig:device}
\end{figure}

For a tilted plane the damping-like torque has a component along the easy axis of the octupole. Reversal
starts when this component exceeds the anisotropy torque, which gives the threshold
\begin{equation}
  \jc(\theta) = \frac{\jzero}{\sin\theta}, \qquad \jzero = \frac{2 e \Ku d}{\hbar\,\tsh},
  \label{eq:threshold}
\end{equation}
where $\Ku$ is the anisotropy energy density and $\tsh$ the spin Hall angle of W (derivation in
Sec.~\ref{S-sec:derivation} of the Supplemental Material~\cite{sm}). The threshold diverges at
$\theta = 0$, so that without tilt no current reverses the moment.

\section{Results}

Figure~\ref{fig:dynamics} shows the response to a current pulse of $6\times10^{11}\un{A/m^2}$. For
$\theta = 0$ the moment tilts by a few degrees and returns. For every tilted film it reverses, and the
reversal is faster for larger tilt. The switching time falls from $0.60\un{ns}$ at $15^\circ$ to
$0.22\un{ns}$ at $45^\circ$. The final state is the same in all runs and is set by the sign of the current,
which is the property a memory needs.

\begin{figure}
  \centering
  \includegraphics[width=0.9\columnwidth]{example-image-b}
  \caption{Out-of-plane component $\mz$ of the octupole moment against time for four tilt angles, at
  $j = 6\times10^{11}\un{A/m^2}$. The pulse starts at $t = 0$.}
  \label{fig:dynamics}
\end{figure}

Figure~\ref{fig:threshold} compares the threshold found in the simulations with Eq.~(\ref{eq:threshold}).
The two agree within $4\%$ between $5^\circ$ and $45^\circ$, with no adjustable parameter. A tilt of
$30^\circ$ costs twice the current of a film tilted by $90^\circ$, and a tilt of $10^\circ$ costs six times
that current. The threshold moves by less than $5\%$ when the damping, the mesh and the rise time of the
pulse are varied over the ranges given in Sec.~\ref{S-sec:robust}.

\begin{figure}
  \centering
  \includegraphics[width=0.9\columnwidth]{example-image-c}
  \caption{Threshold current density $\jc$ against tilt angle $\theta$. Circles show the simulation, and the
  line Eq.~(\ref{eq:threshold}) with the parameters of the simulation.}
  \label{fig:threshold}
\end{figure}

A memory operates near room temperature. At $300\un{K}$ thermal fluctuations broaden the transition from
no reversal to reversal over about $8\%$ of the threshold current, and they do not shift its centre
(Fig.~\ref{S-fig:thermal}). The zero-temperature threshold of Eq.~(\ref{eq:threshold}) therefore remains
the design value.

\section{Conclusion}

In summary, a tilt of the kagome plane makes the current-induced reversal of Mn$_3$Sn deterministic
without a magnetic field. The price is a threshold current that grows inversely with the sine of the tilt
angle, which a simple torque balance predicts without free parameters. Miscut substrates that give tilts
of $30^\circ$ or more keep this price within a factor of two.

\paragraph{Data availability.}
The input files, the raw output of all 212 simulation runs and the scripts that produce every figure of
the paper and of its Supplemental Material are deposited at Zenodo, doi:10.5281/zenodo.0000000
(placeholder).

\begin{thebibliography}{9}
\bibitem{A} [Placeholder: first report of current-induced switching of Mn$_3$Sn with an in-plane field.]
\bibitem{B} [Placeholder: the mechanism of octupole reversal by spin-orbit torque.]
\bibitem{C} [Placeholder: growth of Mn$_3$Sn films with a tilted kagome plane on miscut substrates.]
\bibitem{sm} See Supplemental Material at [URL will be inserted by publisher] for the derivation of
  Eq.~(\ref{eq:threshold}), the robustness of the threshold against numerical and material parameters,
  and the effect of temperature.
\end{thebibliography}
\end{document}
